\documentclass[12pt]{article}
\usepackage[utf8]{inputenc}
\usepackage[T1]{fontenc}
\usepackage{hyperref}
\title{Longto Language Resources for K-12 Curriculum Development}
\author{Compiled by Claude Sonnet 5.5 (claude-sonnet-5-5) for Andr\'as Kornai}
\date{}
\begin{document}
\maketitle

\section{Introduction}
Longto (ISO 639-3 \texttt{wok}; also called Voko, Boko, Gobeyo, Longa, Lonto and Woko) is a Savanna language of the Duru branch. It is spoken in the Poli Subdivision of Faro Department, northern Cameroon. The figure of about 2,400 speakers comes from a 1982 SIL count. Fulfulde is the local lingua franca. Joshua Project reports that young people are increasingly using French as a second language instead of Fulfulde. Wikitongues lists the language as unwritten. Very little material exists, and no school materials were found in the searches.

\section{Primary Educational Resources}
No primers, textbooks or school programmes in Longto were found. Joshua Project mentions a Longto--French dictionary that is hard to find, but no URL for the dictionary itself appeared in search results.

\section{Linguistic and Reference Materials}
\begin{itemize}
\item Wikipedia, Voko language (overview, classification, names): \url{https://en.wikipedia.org/wiki/Voko_language}
\item Glottolog bibliographic entry, Scott Westermann and Annette Westermann 2001: \url{https://glottolog.org/resource/reference/id/310671}
\end{itemize}

\section{Religious and Community Resources}
\begin{itemize}
\item Joshua Project, Longa in Cameroon (people-group profile): \url{https://joshuaproject.net/people_groups/15411}
\end{itemize}
No Longto Scripture or community organisation was found.

\section{Digital and Audio Resources}
\begin{itemize}
\item Wikitongues language page for \texttt{wok}: \url{https://wikitongues.org/languages/wok/}
\item UNESCO World Atlas of Languages entry for Longto: \url{https://ar.wal.unesco.org/en/countries/cameroon/languages/longto}
\end{itemize}
No audio corpus was found.

\section{Curriculum Development Context}
Longto is not known to be taught in schools. Documentation is minimal. A curriculum would first need a documented orthography, a lexicon and community-recorded texts. The UNESCO entry lists use of the language as decreasing, and the sources differ on its vitality: Wikitongues gives EGIDS 6a. Work would also have to take account of Fulfulde and French as the languages of wider communication.
\end{document}